%% ==============================================================
%% КОНФИГУРАЦИЯ — Вариант 1
%% Геометрия A, тип 1 (Эллипсоидальный)
%% Сгенерировано автоматически — не редактировать вручную
%% ==============================================================

% --- Информация о студенте ---
\newcommand{\varNumber}{1}
\newcommand{\varStudent}{Востриков~Я.В.}
\newcommand{\varGroup}{ЗНБ21-02Б}
\newcommand{\varSupervisor}{Башмур~К.А.}
\newcommand{\varStudentID}{}

% --- Геометрия подшипника (геометрия A) ---
\newcommand{\varR}{0.035}
\newcommand{\varRmm}{35}
\newcommand{\varC}{5e-05}
\newcommand{\varCmkm}{50}
\newcommand{\varL}{0.056}
\newcommand{\varLmm}{56}
\newcommand{\varLD}{0.8}
\newcommand{\varGeomLetter}{A}

% --- Тип микрорельефа (тип 1) ---
\newcommand{\varDimpleType}{1}
\newcommand{\varDimpleTypeName}{Эллипсоидальный}
\newcommand{\varDimpleTypeGenitive}{эллипсоидальными}

% --- Параметры микрорельефа ---
\newcommand{\varDimpleA}{0.00241}
\newcommand{\varDimpleAmm}{2.41}
\newcommand{\varDimpleB}{0.002214}
\newcommand{\varDimpleBmm}{2.214}
\newcommand{\varDimpleHp}{1e-05}
\newcommand{\varDimpleHpmkm}{10}
\newcommand{\varDimpleHpRatio}{0.2}

% --- Формула зазора ---
\newcommand{\varDimpleFormula}{
    \Delta h = h_p \sqrt{1 - \left(\frac{\Delta\varphi}{b}\right)^2 - \left(\frac{\Delta Z}{a}\right)^2}
}

% --- Общие параметры ---
\newcommand{\varN}{2980}
\newcommand{\varEta}{0.01105}
\newcommand{\varNphi}{8}
\newcommand{\varNZ}{11}
\newcommand{\varPhiRange}{$90^\circ \div 270^\circ$}
